\documentclass[11pt,a4paper]{article}
\usepackage[utf8]{inputenc}
\usepackage[T1]{fontenc}
\usepackage{lmodern}
\usepackage[italian]{babel}
\usepackage[a4paper,margin=1.9cm,top=2.1cm,bottom=2cm,headheight=14pt]{geometry}
\usepackage{amsmath,booktabs,array,graphicx,float,caption}
\usepackage[expansion=false]{microtype}
\usepackage{pgfplots}
\pgfplotsset{compat=1.18}
\usepackage{xcolor,fancyhdr,xurl}
\usepackage[hidelinks]{hyperref}
\definecolor{navy}{HTML}{193B53}
\definecolor{llm}{HTML}{087F8C}
\definecolor{mqt}{HTML}{CC6B22}
\definecolor{muted}{HTML}{64748B}
\definecolor{lightgrid}{HTML}{E2E8F0}
\pgfplotsset{every axis/.append style={font=\footnotesize,axis line style={muted},
tick style={muted},grid=major,grid style={lightgrid},legend style={draw=none},
scaled ticks=false,/pgf/number format/use comma,/pgf/number format/1000 sep={\,}}}
\captionsetup{font=small,labelfont=bf}
\pagestyle{fancy}\fancyhf{}
\fancyhead[L]{\small\color{muted}Ulteriore test indipendente · QASMBench}
\fancyhead[R]{\small\color{muted}30 settembre 2026}
\fancyfoot[C]{\small\thepage\ / \pageref{ultima}}
\renewcommand{\headrulewidth}{.3pt}
\setlength{\parindent}{0pt}
\setlength{\parskip}{5pt}
\setlength{\emergencystretch}{3em}
\renewcommand{\arraystretch}{1.2}
\setcounter{secnumdepth}{1}
\widowpenalty=10000\clubpenalty=10000
\hypersetup{pdftitle={LLM + RAG e MQT Predictor: ulteriore test indipendente QASMBench},
pdfauthor={Progetto Tesi MQT},pdfsubject={Analisi dei registri del 30 settembre 2026}}
\begin{document}

\begin{center}
{\small\color{muted} VALUTAZIONE DEI SISTEMI DI COMPILAZIONE QUANTISTICA}\\[8pt]
{\LARGE\bfseries\color{navy} LLM + RAG e MQT Predictor}\\[5pt]
{\Large Confronto su QASMBench}\\[9pt]
{\large Ulteriore test indipendente}\\[7pt]
{\small Esecuzioni del 30 settembre 2026 · 50 circuiti · 2 sistemi}
\end{center}
\vspace{8pt}
\section{Risultati principali}
Il test usa circuiti provenienti da QASMBench, diversi dal corpus MQT Bench del progetto.
Sono stati selezionati 30 circuiti piccoli, 15 medi e 5 grandi. Questo documento analizza
i registri conservati; non avvia nuove compilazioni e non modifica le configurazioni valutate.

\begin{center}
\small
\setlength{\tabcolsep}{5pt}
\begin{tabular}{lrr}
\toprule
Indicatore & LLM + RAG & MQT Predictor\\
\midrule
Compilazioni riuscite / circuiti previsti & 50/50 (100\%) & 48/50 (96\%)\\
Score almeno 0,8 / circuiti previsti & 41/50 (82\%) & 43/50 (86\%)\\
Score medio: stessi 48 circuiti & 0,8839 & 0,9168\\
Tempo totale medio: stessi 48 circuiti & 26,76 s & 17,94 s\\
Compilazione interna media: stessi 48 & 0,95 s & 13,50 s\\
\bottomrule
\end{tabular}\end{center}

\paragraph{Copertura.}
LLM + RAG completa tutti i circuiti. MQT Predictor termina entro il limite su 48 circuiti.
I due timeout, su \texttt{gcm\_n13} e \texttt{qft\_n63}, sono riportati separatamente:
non ricevono uno score pari a zero.
\paragraph{Qualità.}
Nei 48 successi comuni, MQT Predictor ottiene uno score medio maggiore di 0,0329. LLM + RAG ottiene lo score maggiore in 22 casi e MQT in 26; non ci sono parità. Le differenze più ampie a favore di MQT si concentrano in alcuni circuiti medi.
\paragraph{Tempi.}
La compilazione interna scelta dall'LLM è più rapida in tutti i 48 successi comuni.
Il tempo necessario a preparare e ottenere la decisione LLM cambia però il risultato
complessivo: MQT è più rapido in 45 di questi 48 casi.
\paragraph{Lettura del confronto.}
Lo score è una fedeltà \emph{stimata} sui Target sintetici, non una misura su hardware quantistico.
Il test confronta i due sistemi completi, con strategie di compilazione diverse.
Non permette di attribuire il risultato al solo selettore del dispositivo.

\clearpage\section{Circuiti e impostazioni}

La selezione è fissata nel manifest della campagna, prima di queste analisi. È una
selezione ragionata, non un campione casuale dell'intera raccolta QASMBench.
Non sono state trovate copie identiche byte per byte nel corpus MQT controllato.
Questo controllo non esclude equivalenze algoritmiche o la presenza di circuiti simili
nell'addestramento originario dell'LLM.

\begin{center}
\small
\setlength{\tabcolsep}{5pt}
\begin{tabular}{lrrr}
\toprule
Fascia & Circuiti & Quota & Qubit osservati\\
\midrule
Piccoli & 30 & 60\% & 2--10\\
Medi & 15 & 30\% & 11--27\\
Grandi & 5 & 10\% & 28, 63, 98, 111, 140\\
\bottomrule
\end{tabular}\end{center}

\begin{center}
\small
\setlength{\tabcolsep}{5pt}
\begin{tabular}{p{4.2cm}p{11.1cm}}
\toprule
Voce & Impostazione\\
\midrule
LLM & Qwen3.5-4B, pesi Q8\_0, temperatura 0\\
Generazione & Seed 20260913; massimo 4096 token in uscita\\
Altri parametri & top\_p 0,95; top\_k 40; min\_p 0; cache attiva\\
Contesto & 60\,000 token richiesti; 60\,160 dichiarati dal server\\
RAG & 5 esempi; distanza Manhattan; 49 caratteristiche\\
Dataset RAG & 396 circuiti train distinti; formato del prompt TOON\\
Decisione & Contratto 4.0.0; al massimo 3 risposte completate\\
MQT & mqt.predictor 2.4.0; selettore supervisionato + politiche RL\\
Training set MQT & 384 dei 396 campioni previsti; 12 esclusi\\
Raccolta del Training set & 1853/1878 compilazioni riuscite; limiti adattivi 100/300 s\\
Compilazione del Test & Una per circuito e metodo; limite del processo: 100 s\\
Semi della compilazione & Qiskit: 0; campionamento MQT: 0\\
Ambiente & Python 3.12.13; WSL2 x86\_64; 12 CPU logiche rilevate\\
\bottomrule
\end{tabular}\end{center}

\paragraph{Misura della qualità.}
Per ciascuna istruzione non di barriera del circuito compilato si usa l'errore
registrato nel Target. La metrica riproduce \texttt{expected\_fidelity} di MQT Predictor 2.4.0

\clearpage\section{Riuscita e soglia di qualità}

\begin{figure}[H]\centering
\begin{tikzpicture}\begin{axis}[width=\textwidth,height=6.1cm,ybar stacked,bar width=34pt,ymin=0,ymax=54,ytick={0,10,20,30,40,50},ylabel={Numero di circuiti},xtick={1,2},xticklabels={LLM + RAG,MQT Predictor},xmin=.4,xmax=2.6,legend style={at={(.5,-.18)},anchor=north,legend columns=3,font=\small}]
\addplot[ybar,fill=llm,draw=llm] coordinates {(1,41) (2,43)};
\addlegendentry{Score almeno 0,8}
\addplot[ybar,fill=mqt,draw=mqt] coordinates {(1,9) (2,5)};
\addlegendentry{Score sotto 0,8}
\addplot[ybar,fill=muted,draw=muted] coordinates {(1,0) (2,2)};
\addlegendentry{Timeout}
\end{axis}\end{tikzpicture}
\caption{Tutti i 50 circuiti per sistema. I timeout sono distinti dai successi sotto soglia. La soglia 0,8 riprende il report di riferimento; non è una garanzia di accuratezza fisica.}
\end{figure}

\begin{center}
\small
\setlength{\tabcolsep}{5pt}
\begin{tabular}{lrrrr}
\toprule
Fascia & Successi LLM & Successi MQT & $S\geq0{,}8$ LLM & $S\geq0{,}8$ MQT\\
\midrule
Totale & 50/50 & 48/50 & 41/50 & 43/50\\
Piccoli & 30/30 & 30/30 & 30/30 & 29/30\\
Medi & 15/15 & 14/15 & 11/15 & 14/15\\
Grandi & 5/5 & 4/5 & 0/5 & 0/5\\
\bottomrule
\end{tabular}\end{center}

Nei piccoli circuiti LLM + RAG supera la soglia in tutti i casi. Nei medi,
MQT la supera in tutti i 14 casi completati, mentre LLM + RAG la supera in 11 casi su 15.
Nessuno dei due sistemi raggiunge 0,8 sui grandi circuiti completati.
\subsection*{I due timeout di MQT Predictor}
Il limite di 100 secondi riguarda l'intero processo di compilazione: include avvio,
importazioni e selezione del dispositivo. La durata registrata supera leggermente
il limite per l'arresto e la raccolta dell'esito.

\begin{center}
\small
\setlength{\tabcolsep}{5pt}
\begin{tabular}{lrrrrr}
\toprule
Circuito & Qubit & Target MQT & $S$ LLM & Totale LLM & Totale MQT\\
\midrule
gcm\_n13 & 13 & Q56 & 0,533621 & 27,60 s & 100,17 s\\
qft\_n63 & 63 & H156 & 0,000076 & 22,01 s & 100,18 s\\
\bottomrule
\end{tabular}\end{center}

I Target dei due timeout sono ricavati da \texttt{selection.json}, scritto prima della
compilazione: Q56 indica Quantinuum H2-56 e H156 IBM Heron 156.
La scelta era quindi già avvenuta. Non è disponibile uno score MQT da confrontare.
LLM + RAG completa \texttt{qft\_n63}, ma con uno score molto basso: riuscita tecnica
e qualità del risultato restano due misure distinte.

\clearpage\section{Qualità sui 48 successi comuni}

\begin{figure}[H]\centering
\begin{minipage}[t]{.48\textwidth}\centering\begin{tikzpicture}\begin{axis}[width=\linewidth,height=6.8cm,xmin=0,xmax=1,ymin=0,ymax=1,xlabel={Score MQT},ylabel={Score LLM + RAG},legend style={at={(.5,-.25)},anchor=north,legend columns=3,font=\scriptsize}]
\addplot[gray,dashed,no marks,forget plot] coordinates {(0,0) (1,1)};
\addplot[only marks,mark=*,color=llm,mark size=2pt] coordinates {(0.9933037999,0.99930613) (0.9847393348,0.9990727123) (0.9837216237,0.9979257622) (0.9875104411,0.9990727123) (0.9949096128,0.9965837357) (0.9827085171,0.9890425841) (0.9899469686,0.9893315188) (0.9850524927,0.9852437327) (0.9913896853,0.9879867544) (0.9770013741,0.986691458) (0.9965093604,0.9992292817) (0.9880639018,0.991055686) (0.9488232061,0.9055163799) (0.9898451376,0.9799502249) (0.9862858714,0.9939425554) (0.9916295845,0.991786503) (0.9931192634,0.9812297456) (0.9880639018,0.9813525305) (0.9774435979,0.9899937206) (0.9607174974,0.9198463434) (0.7261481348,0.8015117741) (0.9882121247,0.9891612781) (0.9740750274,0.9765914234) (0.9896966697,0.9795571896) (0.9822972636,0.971707646) (0.9901421401,0.9185065432) (0.9427187379,0.8958629829) (0.9727598201,0.9016650371) (0.9832713192,0.9499308505) (0.9730516962,0.9157132407)};
\addlegendentry{Piccoli}
\addplot[only marks,mark=triangle*,color=mqt,mark size=2pt] coordinates {(0.9607174974,0.7095805071) (0.9632015121,0.8467850023) (0.9519626937,0.9593308041) (0.9444192698,0.7111796696) (0.9462616675,0.8129998391) (0.9327526847,0.9376907179) (0.9689701802,0.7719882283) (0.8781117552,0.8212656134) (0.8972984622,0.8248486851) (0.9339256194,0.9344861513) (0.9253337216,0.9315450511) (0.9475441855,0.8716948434) (0.9344051438,0.9230263603) (0.8807706766,0.8819066477)};
\addlegendentry{Medi}
\addplot[only marks,mark=square*,color=navy,mark size=2pt] coordinates {(0.7988808034,0.728487822) (0.6527318587,0.5221888556) (0.0030146815,0.0114679021) (0.2725048164,0.2620161724)};
\addlegendentry{Grandi}
\end{axis}\end{tikzpicture}
\end{minipage}\hfill\begin{minipage}[t]{.48\textwidth}\centering\begin{tikzpicture}\begin{axis}[width=\linewidth,height=6.8cm,xmin=0,xmax=1,ymin=0,ymax=100,xlabel={Score},ylabel={Circuiti con score non superiore (\%)},legend style={at={(.5,-.25)},anchor=north,legend columns=1,font=\scriptsize}]
\addplot[const plot,thick,color=llm] coordinates {(0.0114679021,2.083333333) (0.2620161724,4.166666667) (0.5221888556,6.25) (0.7095805071,8.333333333) (0.7111796696,10.41666667) (0.728487822,12.5) (0.7719882283,14.58333333) (0.8015117741,16.66666667) (0.8129998391,18.75) (0.8212656134,20.83333333) (0.8248486851,22.91666667) (0.8467850023,25) (0.8716948434,27.08333333) (0.8819066477,29.16666667) (0.8958629829,31.25) (0.9016650371,33.33333333) (0.9055163799,35.41666667) (0.9157132407,37.5) (0.9185065432,39.58333333) (0.9198463434,41.66666667) (0.9230263603,43.75) (0.9315450511,45.83333333) (0.9344861513,47.91666667) (0.9376907179,50) (0.9499308505,52.08333333) (0.9593308041,54.16666667) (0.971707646,56.25) (0.9765914234,58.33333333) (0.9795571896,60.41666667) (0.9799502249,62.5) (0.9812297456,64.58333333) (0.9813525305,66.66666667) (0.9852437327,68.75) (0.986691458,70.83333333) (0.9879867544,72.91666667) (0.9890425841,75) (0.9891612781,77.08333333) (0.9893315188,79.16666667) (0.9899937206,81.25) (0.991055686,83.33333333) (0.991786503,85.41666667) (0.9939425554,87.5) (0.9965837357,89.58333333) (0.9979257622,91.66666667) (0.9990727123,93.75) (0.9990727123,95.83333333) (0.9992292817,97.91666667) (0.99930613,100)};
\addlegendentry{LLM + RAG}
\addplot[const plot,thick,color=mqt] coordinates {(0.0030146815,2.083333333) (0.2725048164,4.166666667) (0.6527318587,6.25) (0.7261481348,8.333333333) (0.7988808034,10.41666667) (0.8781117552,12.5) (0.8807706766,14.58333333) (0.8972984622,16.66666667) (0.9253337216,18.75) (0.9327526847,20.83333333) (0.9339256194,22.91666667) (0.9344051438,25) (0.9427187379,27.08333333) (0.9444192698,29.16666667) (0.9462616675,31.25) (0.9475441855,33.33333333) (0.9488232061,35.41666667) (0.9519626937,37.5) (0.9607174974,39.58333333) (0.9607174974,41.66666667) (0.9632015121,43.75) (0.9689701802,45.83333333) (0.9727598201,47.91666667) (0.9730516962,50) (0.9740750274,52.08333333) (0.9770013741,54.16666667) (0.9774435979,56.25) (0.9822972636,58.33333333) (0.9827085171,60.41666667) (0.9832713192,62.5) (0.9837216237,64.58333333) (0.9847393348,66.66666667) (0.9850524927,68.75) (0.9862858714,70.83333333) (0.9875104411,72.91666667) (0.9880639018,75) (0.9880639018,77.08333333) (0.9882121247,79.16666667) (0.9896966697,81.25) (0.9898451376,83.33333333) (0.9899469686,85.41666667) (0.9901421401,87.5) (0.9913896853,89.58333333) (0.9916295845,91.66666667) (0.9931192634,93.75) (0.9933037999,95.83333333) (0.9949096128,97.91666667) (0.9965093604,100)};
\addlegendentry{MQT Predictor}
\end{axis}\end{tikzpicture}
\end{minipage}\caption{A sinistra ogni punto è un circuito: sopra la diagonale prevale LLM + RAG. A destra la distribuzione cumulativa sui medesimi 48 circuiti: a parità di score, una quota più bassa indica meno circuiti con qualità bassa.}
\end{figure}

\begin{center}
\small
\setlength{\tabcolsep}{5pt}
\begin{tabular}{lrrrrr}
\toprule
Popolazione & Sistema & N & Media & Mediana & Minimo\\
\midrule
Successi comuni & LLM + RAG & 48 & 0,8839 & 0,9438 & 0,011468\\
Successi comuni & MQT Predictor & 48 & 0,9168 & 0,9736 & 0,003015\\
Tutti i successi & LLM + RAG & 50 & 0,8592 & 0,9361 & 0,000076\\
Tutti i successi & MQT Predictor & 48 & 0,9168 & 0,9736 & 0,003015\\
\bottomrule
\end{tabular}\end{center}

La prima coppia di righe è il confronto principale: entrambi gli score esistono
per gli stessi circuiti. Le medie su tutti i successi hanno denominatori diversi
e servono soltanto a descrivere ciascun sistema.

\clearpage\section{Dove cambiano maggiormente gli score}

\begin{figure}[H]\centering
\begin{tikzpicture}\begin{axis}[width=\textwidth,height=6.4cm,xmin=0,xmax=51,ymin=-.28,ymax=.1,xtick={1,10,20,30,40,50},xlabel={Indice del circuito nel manifest},ylabel={Differenza di score},legend style={at={(.5,-.22)},anchor=north,legend columns=3}]
\addplot[ybar,bar width=3pt,fill=llm,draw=llm] coordinates {(1,0.0060023301) (2,0.0143333775) (3,0.0142041385) (4,0.0115622712) (5,0.0016741229) (6,0.006334067) (7,-0.0006154498) (8,0.00019124) (9,-0.0034029309) (10,0.0096900839) (11,0.0027199213) (12,0.0029917842) (13,-0.0433068262) (14,-0.0098949127) (15,0.007656684) (16,0.0001569185) (17,-0.0118895178) (18,-0.0067113713) (19,0.0125501227) (20,-0.040871154) (21,0.0753636393) (22,0.0009491534) (23,0.002516396) (24,-0.0101394801) (25,-0.0105896176) (26,-0.0716355969) (27,-0.046855755) (28,-0.071094783) (29,-0.0333404687) (30,-0.0573384555)};
\addlegendentry{Piccoli}
\addplot[ybar,bar width=3pt,fill=mqt,draw=mqt] coordinates {(31,-0.2511369903) (32,-0.1164165098) (33,0.0073681104) (34,-0.2332396002) (35,-0.1332618284) (36,0.0049380332) (37,-0.1969819519) (38,-0.0568461418) (40,-0.0724497771) (41,0.0005605319) (42,0.0062113295) (43,-0.0758493421) (44,-0.0113787835) (45,0.0011359711)};
\addlegendentry{Medi}
\addplot[ybar,bar width=3pt,fill=navy,draw=navy] coordinates {(46,-0.0703929814) (48,-0.1305430031) (49,0.0084532206) (50,-0.010488644)};
\addlegendentry{Grandi}
\addplot[black,thin,no marks,forget plot] coordinates {(0,0) (51,0)};
\end{axis}\end{tikzpicture}
\caption{Differenze sui singoli circuiti. Valori positivi favoriscono LLM + RAG. Gli indici 39 e 47 non hanno una barra: sono i timeout MQT. L'appendice associa gli indici ai nomi.}
\end{figure}

\begin{center}
\small
\setlength{\tabcolsep}{5pt}
\begin{tabular}{lrrrr}
\toprule
Circuito & LLM & MQT & $\Delta$ & Target L/M\\
\midrule
sat\_n11 & 0,7096 & 0,9607 & -0,2511 & Q56/Q56\\
multiplier\_n15 & 0,7112 & 0,9444 & -0,2332 & Q56/Q56\\
qft\_n18 & 0,7720 & 0,9690 & -0,1970 & Q56/Q56\\
dnn\_n16 & 0,8130 & 0,9463 & -0,1333 & Q56/Q56\\
basis\_trotter\_n4 & 0,8015 & 0,7261 & +0,0754 & Q56/Q56\\
iswap\_n2 & 0,9991 & 0,9847 & +0,0143 & F127/F127\\
quantumwalks\_n2 & 0,9979 & 0,9837 & +0,0142 & F127/F127\\
vqe\_n4 & 0,9900 & 0,9774 & +0,0126 & F127/Q56\\
\bottomrule
\end{tabular}\end{center}

La tabella mostra i quattro maggiori scarti in ciascuna direzione,
senza selezionare i casi in base alla narrativa. I nomi dei dispositivi sono
abbreviati: F127 = IBM Falcon 127, H133/H156 = IBM Heron 133/156,
Q56 = Quantinuum H2-56.


Il maggiore vantaggio MQT compare su \texttt{sat\_n11}: lo scarto è 0,2511. Il maggiore vantaggio LLM + RAG compare su \texttt{basis\_trotter\_n4}: 0,0754.

La sola scelta dello stesso dispositivo non rende uguali i sistemi.
LLM + RAG sceglie anche una configurazione Qiskit dal catalogo; MQT usa una politica
RL per costruire la sequenza dei passi. Differenze di dispositivo e di compilazione
agiscono insieme. Questi dati non isolano una causa unica del vantaggio osservato.

\clearpage\section{Costo delle decisioni LLM e correzioni}

\begin{center}
\small
\setlength{\tabcolsep}{5pt}
\begin{tabular}{lrrr}
\toprule
Misura & Totale & Media per circuito & Mediana\\
\midrule
Token in ingresso & 903\,142 & 18\,063 & 12\,067\\
Token in uscita & 12\,529 & 251 & 158\\
Token complessivi & 915\,671 & 18\,313 & 12\,231\\
Chiamate LLM & 78,00 & 1,56 & 1,00\\
Chiamate aggiuntive & 28,00 & 0,56 & 0,00\\
Tempo RAG (s) & 179,37 & 3,59 & 3,63\\
Tempo risposte LLM (s) & 1\,014,57 & 20,29 & 19,65\\
\bottomrule
\end{tabular}\end{center}

\begin{figure}[H]\centering
\begin{minipage}[t]{.48\textwidth}\centering\begin{tikzpicture}\begin{axis}[width=\linewidth,height=5.6cm,ymin=0,ymax=40,xmin=.5,xmax=3.5,xtick={1,2,3},xlabel={Chiamate per circuito},ylabel={Numero di circuiti},bar width=24pt]
\addplot[ybar,fill=llm,draw=llm,nodes near coords] coordinates {(1,35) (2,2) (3,13)};
\end{axis}\end{tikzpicture}
\end{minipage}\hfill\begin{minipage}[t]{.48\textwidth}\centering\begin{tikzpicture}\begin{axis}[width=\linewidth,height=5.6cm,ymin=0,ymax=40,xmin=.5,xmax=3.5,xtick={1,2,3},xlabel={Chiamate per circuito},ylabel={Token medi (migliaia)},bar width=24pt]
\addplot[ybar,fill=mqt,draw=mqt,nodes near coords,point meta=y] coordinates {(1,11.67048571) (2,23.375) (3,35.41953846)};
\end{axis}\end{tikzpicture}
\end{minipage}\caption{Le correzioni aumentano il numero di richieste e il volume di token. Le colonne di destra sono medie su 35, 2 e 13 circuiti, rispettivamente. Non misurano l'effetto causale delle correzioni sulla qualità.}
\end{figure}

\clearpage\section{Scelte dei sistemi e limiti del confronto}

\begin{center}
\small
\setlength{\tabcolsep}{5pt}
\begin{tabular}{lrrr}
\toprule
Dispositivo & LLM (50 successi) & MQT (48 successi) & MQT timeout\\
\midrule
IBM Falcon 27 & 0 & 0 & 0\\
IBM Falcon 127 & 12 & 8 & 0\\
IBM Heron 133 & 0 & 4 & 0\\
IBM Heron 156 & 10 & 1 & 1\\
Quantinuum H2-56 & 28 & 35 & 1\\
\bottomrule
\end{tabular}\end{center}

LLM + RAG sceglie 24 volte \texttt{o2\_default\_default}, 24 volte
\texttt{o3\_default\_default} e 2 volte \texttt{o2\_dense\_sabre}.
MQT non sceglie una di queste configurazioni: esegue i passi della politica RL.
Il suo selettore conserva quattro classi di dispositivo (Falcon 127, Heron 133,
Heron 156 e H2-56). Falcon 27 è disponibile nel catalogo, ma non è una classe
appresa da questo selettore. Le frequenze riflettono anche capacità e compatibilità
dei dispositivi con i circuiti.

\paragraph{Ambito della conclusione.}
MQT ottiene una qualità media maggiore sui 48 circuiti confrontabili; LLM + RAG
offre copertura completa entro il limite impostato. La rapidità della compilazione
Qiskit non basta a rendere più rapido l'intero sistema LLM.
Sono tre risultati distinti: qualità, riuscita e tempo.

\paragraph{Vincoli dei modelli.}
Il confronto riguarda gli artefatti realmente usati. Il Training set del selettore
MQT comprende 384 campioni, non tutti i 396 previsti. La raccolta che lo ha prodotto
ha 1853 compilazioni riuscite su 1878 e ha usato limiti adattivi di 100/300 secondi.
Questi limiti di addestramento sono diversi dai 100 secondi del Test. Il documento
non estende i risultati a un selettore ritrainato su dati diversi.


\clearpage\section{Appendice: risultati per circuito}

A = LLM + RAG; B = MQT Predictor. $S$ è lo score; $t$ il tempo totale
in secondi, comprensivo della scelta. P/M/G indicano piccolo/medio/grande.
TO indica un timeout: il tempo dell'arresto è osservato, lo score è assente.
Il prefisso \texttt{qasmbench\_} è omesso dai nomi.

\begingroup\footnotesize\renewcommand{\arraystretch}{1.45}

\begin{center}
\setlength{\tabcolsep}{5pt}
\begin{tabular}{rlrrrrrrrr}
\toprule
N. & Circuito & Qubit & F. & $S_A$ & $S_B$ & $t_A$ & $t_B$ & Ch. & \\
\midrule
1 & deutsch\_n2 & 2 & P & 0,99931 & 0,99330 & 22,4 & 14,8 & 1 & \\
2 & iswap\_n2 & 2 & P & 0,99907 & 0,98474 & 22,8 & 14,3 & 1 & \\
3 & quantumwalks\_n2 & 2 & P & 0,99793 & 0,98372 & 22,6 & 12,1 & 1 & \\
4 & grover\_n2 & 2 & P & 0,99907 & 0,98751 & 29,5 & 14,0 & 3 & *\\
5 & dnn\_n2 & 2 & P & 0,99658 & 0,99491 & 27,2 & 15,8 & 2 & \\
6 & qaoa\_n3 & 3 & P & 0,98904 & 0,98271 & 23,3 & 16,7 & 1 & \\
7 & toffoli\_n3 & 3 & P & 0,98933 & 0,98995 & 29,3 & 14,5 & 3 & *\\
8 & fredkin\_n3 & 3 & P & 0,98524 & 0,98505 & 22,8 & 14,7 & 1 & \\
9 & wstate\_n3 & 3 & P & 0,98799 & 0,99139 & 22,7 & 14,3 & 1 & \\
10 & basis\_change\_n3 & 3 & P & 0,98669 & 0,97700 & 23,4 & 17,5 & 1 & \\
11 & qrng\_n4 & 4 & P & 0,99923 & 0,99651 & 24,1 & 14,5 & 1 & \\
12 & cat\_state\_n4 & 4 & P & 0,99106 & 0,98806 & 23,0 & 14,0 & 1 & \\
13 & adder\_n10 & 10 & P & 0,90552 & 0,94882 & 27,1 & 26,1 & 1 & \\
14 & adder\_n4 & 4 & P & 0,97995 & 0,98985 & 22,8 & 14,5 & 1 & \\
15 & hs4\_n4 & 4 & P & 0,99394 & 0,98629 & 29,3 & 16,0 & 3 & *\\
16 & bell\_n4 & 4 & P & 0,99179 & 0,99163 & 27,0 & 16,3 & 2 & \\
17 & qft\_n4 & 4 & P & 0,98123 & 0,99312 & 29,3 & 14,3 & 3 & *\\
18 & variational\_n4 & 4 & P & 0,98135 & 0,98806 & 31,4 & 14,5 & 3 & *\\
19 & vqe\_n4 & 4 & P & 0,98999 & 0,97744 & 29,8 & 14,4 & 3 & *\\
20 & sat\_n7 & 7 & P & 0,91985 & 0,96072 & 30,3 & 15,6 & 1 & \\
21 & basis\_trotter\_n4 & 4 & P & 0,80151 & 0,72615 & 30,2 & 54,7 & 3 & *\\
22 & lpn\_n5 & 5 & P & 0,98916 & 0,98821 & 23,5 & 15,7 & 1 & \\
23 & qec\_en\_n5 & 5 & P & 0,97659 & 0,97408 & 24,8 & 14,4 & 1 & \\
24 & pea\_n5 & 5 & P & 0,97956 & 0,98970 & 30,0 & 17,8 & 3 & *\\
25 & simon\_n6 & 6 & P & 0,97171 & 0,98230 & 24,8 & 14,4 & 1 & \\
\bottomrule
\end{tabular}\end{center}

\endgroup

Ch. indica il numero di chiamate LLM. L'asterisco segnala una decisione
accettata con fatti non completamente verificati. Gli score sono visualizzati
con cinque decimali: confronti, soglie e differenze sono calcolati dai dieci decimali
conservati nei registri, prima di questa formattazione.

I file \path{dati/circuiti.csv} e \path{dati/coppie.csv} riportano anche
tempi interni e di processo, dispositivi, configurazioni, token, correzioni,
log-score, stati e impronte dei circuiti. I dati mancanti restano vuoti,
senza sostituzione con zeri.

\clearpage\section{Appendice: risultati per circuito (continua)}

A = LLM + RAG; B = MQT Predictor. $S$ è lo score; $t$ il tempo totale
in secondi, comprensivo della scelta. P/M/G indicano piccolo/medio/grande.
TO indica un timeout: il tempo dell'arresto è osservato, lo score è assente.
Il prefisso \texttt{qasmbench\_} è omesso dai nomi.

\begingroup\footnotesize\renewcommand{\arraystretch}{1.45}

\begin{center}
\setlength{\tabcolsep}{5pt}
\begin{tabular}{rlrrrrrrrr}
\toprule
N. & Circuito & Qubit & F. & $S_A$ & $S_B$ & $t_A$ & $t_B$ & Ch. & \\
\midrule
26 & qaoa\_n6 & 6 & P & 0,91851 & 0,99014 & 24,3 & 17,2 & 1 & \\
27 & hhl\_n7 & 7 & P & 0,89586 & 0,94272 & 25,3 & 21,1 & 1 & \\
28 & dnn\_n8 & 8 & P & 0,90167 & 0,97276 & 25,3 & 18,8 & 1 & \\
29 & qpe\_n9 & 9 & P & 0,94993 & 0,98327 & 33,7 & 16,4 & 3 & *\\
30 & ising\_n10 & 10 & P & 0,91571 & 0,97305 & 30,4 & 18,8 & 3 & *\\
31 & sat\_n11 & 11 & M & 0,70958 & 0,96072 & 27,7 & 24,4 & 1 & \\
32 & multiply\_n13 & 13 & M & 0,84679 & 0,96320 & 30,0 & 15,1 & 3 & *\\
33 & bv\_n14 & 14 & M & 0,95933 & 0,95196 & 26,4 & 18,3 & 1 & \\
34 & multiplier\_n15 & 15 & M & 0,71118 & 0,94442 & 26,3 & 24,3 & 1 & \\
35 & dnn\_n16 & 16 & M & 0,81300 & 0,94626 & 60,2 & 16,4 & 3 & *\\
36 & qec9xz\_n17 & 17 & M & 0,93769 & 0,93275 & 25,2 & 14,9 & 1 & \\
37 & qft\_n18 & 18 & M & 0,77199 & 0,96897 & 25,6 & 19,6 & 1 & \\
38 & bigadder\_n18 & 18 & M & 0,82127 & 0,87811 & 27,6 & 19,7 & 1 & \\
39 & gcm\_n13 & 13 & M & 0,53362 & TO & 27,6 & 100,2 & 1 & \\
40 & qram\_n20 & 20 & M & 0,82485 & 0,89730 & 25,5 & 19,0 & 1 & \\
41 & cat\_state\_n22 & 22 & M & 0,93449 & 0,93393 & 25,1 & 12,5 & 1 & \\
42 & ghz\_state\_n23 & 23 & M & 0,93155 & 0,92533 & 24,9 & 15,1 & 1 & \\
43 & swap\_test\_n25 & 25 & M & 0,87169 & 0,94754 & 25,2 & 33,8 & 1 & \\
44 & ising\_n26 & 26 & M & 0,92303 & 0,93441 & 30,5 & 14,6 & 3 & *\\
45 & wstate\_n27 & 27 & M & 0,88191 & 0,88077 & 25,2 & 22,0 & 1 & \\
46 & adder\_n28 & 28 & G & 0,72849 & 0,79888 & 24,4 & 23,0 & 1 & \\
47 & qft\_n63 & 63 & G & 0,00008 & TO & 22,0 & 100,2 & 1 & \\
48 & ising\_n98 & 98 & G & 0,52219 & 0,65273 & 21,9 & 23,0 & 1 & \\
49 & qugan\_n111 & 111 & G & 0,01147 & 0,00301 & 26,5 & 13,8 & 1 & \\
50 & bv\_n140 & 140 & G & 0,26202 & 0,27250 & 14,4 & 13,5 & 1 & \\
\bottomrule
\end{tabular}\end{center}

\endgroup

Ch. indica il numero di chiamate LLM. L'asterisco segnala una decisione
accettata con fatti non completamente verificati. Gli score sono visualizzati
con cinque decimali: confronti, soglie e differenze sono calcolati dai dieci decimali
conservati nei registri, prima di questa formattazione.

I file \path{dati/circuiti.csv} e \path{dati/coppie.csv} riportano anche
tempi interni e di processo, dispositivi, configurazioni, token, correzioni,
log-score, stati e impronte dei circuiti. I dati mancanti restano vuoti,
senza sostituzione con zeri.

\label{ultima}\end{document}
